\chapter{प्रकाश के वर्णक्रम और प्रकाश का संदेश}\label{ch:g10:light-spectra}

तारे की रोशनी की एक किरण ही उस तारे का इकलौता नमूना है जो कभी हमारे हाथ
आएगा, और वह नमूना बड़ा उदार निकलता है। किरण को उसके \omterm{def:g10:light-spectra:white}{वर्णक्रम} में
फैला दीजिए --- काँच का प्रिज़्म यह कर देता है, बारिश की बूँदों का परदा
भी --- और प्रकाश एक संदेश बन जाता है: पीछे फैली सतत इंद्रधनुषी पट्टी
तारे का ताप बता देती है, और महीन रेखाओं का बारकोड एक-एक तत्व गिनाकर
बताता है कि तारा किस चीज़ का बना है। यह अध्याय संदेश के दोनों हिस्से
पढ़ना सिखाता है --- पहले प्रयोगशाला के लैंपों पर, फिर सूर्य पर।

\section{श्वेत प्रकाश और उसका वर्णक्रम}

\begin{definition}[श्वेत प्रकाश का वर्ण-विक्षेपण]\label{def:g10:light-spectra:white}
जब धूप की एक पतली किरण काँच के प्रिज़्म से गुज़रती है, तो वह लाल से
बैंगनी तक फैली रंगों की एक सतत पट्टी बनकर निकलती है। यही पट्टी उस प्रकाश
का \emph{वर्णक्रम}\index{वर्णक्रम} है। जिस प्रकाश में सभी दृश्य रंग
मौजूद हों --- जैसे धूप या साधारण बल्ब का प्रकाश --- उसे \emph{श्वेत
प्रकाश}\index{श्वेत प्रकाश} कहते हैं। एक ही शुद्ध रंग का प्रकाश, जिसे
कोई प्रिज़्म आगे नहीं बाँट सकता, \emph{एकवर्णी}\index{एकवर्णी प्रकाश}
कहलाता है (जैसे लेज़र की किरण); कई रंगों का मिश्रण
\emph{बहुवर्णी}\index{बहुवर्णी प्रकाश} होता है। जो उपकरण यह विक्षेपण
करके किसी प्रकाश की बनावट पढ़ लेता है, वह
\emph{वर्णक्रमदर्शी}\index{वर्णक्रमदर्शी} है।
\end{definition}

\begin{omfigure}
\begin{tikzpicture}[scale=1.0, font=\small]
  \draw[thick, fill=blue!6] (0,0) -- (1.4,2.6) -- (2.8,0) -- cycle;
  \draw[line width=1.6pt] (-2.4,1.95) -- (0.75,1.4)
    node[pos=0.22, above] {श्वेत प्रकाश};
  \draw[gray] (0.75,1.4) -- (2.05,1.05);
  \draw[red, thick]              (2.05,1.05) -- (5.4,0.30);
  \draw[orange, thick]           (2.05,1.05) -- (5.36,0.08);
  \draw[yellow!75!black, thick]  (2.05,1.05) -- (5.31,-0.14);
  \draw[green!65!black, thick]   (2.05,1.05) -- (5.25,-0.36);
  \draw[blue, thick]             (2.05,1.05) -- (5.17,-0.58);
  \draw[violet, thick]           (2.05,1.05) -- (5.08,-0.80);
  \node[right] at (5.4,0.30)  {लाल (सबसे कम विचलित)};
  \node[right] at (5.08,-0.80) {बैंगनी (सबसे अधिक विचलित)};
\end{tikzpicture}

{\small प्रिज़्म \omterm{def:g10:light-spectra:white}{श्वेत प्रकाश} की किरण को उसके \omterm{def:g10:light-spectra:white}{वर्णक्रम} में
पंखे की तरह फैला देता है; परदे पर यह पंखा एक सतत इंद्रधनुषी पट्टी
रँगता है।}
\end{omfigure}

\begin{remark}[प्रिज़्म और बारिश की बूँदें]\label{rem:g10:light-spectra:rainbow}
प्रिज़्म इसलिए काम करता है कि काँच बैंगनी प्रकाश को लाल से थोड़ा अधिक
मोड़ता है --- ऐसा \emph{क्यों} होता है, यह अपवर्तन के अध्ययन का विषय है,
\cref{ch:g10:refraction} में; यहाँ हम केवल प्रभाव का उपयोग करते हैं। पानी भी यही काम करता
है: गुज़रती बौछार की हर बूँद एक नन्हे प्रिज़्म की तरह काम करती है, श्वेत
धूप लेती है और उसे रंगवार छाँटकर लौटा देती है। इंद्रधनुष आकाश भर में खिंचा
हुआ सूर्य का \omterm{def:g10:light-spectra:white}{वर्णक्रम} ही है।
\end{remark}

\begin{definition}[तरंगदैर्घ्य, किसी विकिरण का पहचान-पत्र]\label{def:g10:light-spectra:wavelength}
हर \omterm{def:g10:light-spectra:white}{एकवर्णी} विकिरण एक ही संख्या से पहचाना जाता है, उसकी
\emph{तरंगदैर्घ्य}\index{तरंगदैर्घ्य} $\lambda$, जिसे नैनोमीटर ($\qty{1}{nm} = \qty{e-9}{m}$) में
मापा जाता है। तरंगदैर्घ्य विकिरण का पहचान-पत्र है: $\lambda$ बता दीजिए और
उस शुद्ध रंग के बारे में कहने को कुछ शेष नहीं रहता। आँख लगभग $\qty{400}{nm}$
(बैंगनी) और $\qty{800}{nm}$ (लाल) के बीच की तरंगदैर्घ्यों पर प्रतिक्रिया करती है
--- यही \emph{दृश्य वर्णक्रम}\index{दृश्य वर्णक्रम} है। दोनों सिरों के
आगे विकिरण चलता रहता है, पर हमें दिखता नहीं: $\qty{400}{nm}$ से नीचे
\emph{पराबैंगनी}\index{पराबैंगनी}, $\qty{800}{nm}$ से ऊपर
\emph{अवरक्त}\index{अवरक्त}। तरंगदैर्घ्य वास्तव में क्या नापती है ---
प्रकाश तरंग की एक लहर की लंबाई --- यह \cref{ch:g10:signals-and-waves} में समझाया गया है; इस
अध्याय में संख्या ही काम की है।
\end{definition}

\begin{example}[तरंगदैर्घ्य पढ़ना]\label{ex:g10:light-spectra:classify}
हरे लेज़र सूचक का उत्सर्जन $\qty{532}{nm}$ पर होता है: दृश्य परिसर के भीतर, हरे
में। कीटाणुनाशक पारा लैंप $\qty{254}{nm}$ पर उत्सर्जन करता है: \omterm{def:g10:light-spectra:wavelength}{पराबैंगनी},
अदृश्य (और आँख के लिए ख़तरनाक ठीक इसीलिए कि अदृश्य है)। टेलीविज़न रिमोट
का डायोड $\qty{940}{nm}$ पर उत्सर्जन करता है: \omterm{def:g10:light-spectra:wavelength}{अवरक्त} --- उसे फ़ोन के
कैमरे की ओर कीजिए, जिसका संवेदक $\qty{800}{nm}$ से थोड़ा आगे तक देख लेता है, और
अदृश्य चमक परदे पर उभर आएगी।
\end{example}

\section{सतत वर्णक्रम और ताप}

\begin{definition}[सतत वर्णक्रम]\label{def:g10:light-spectra:continuous}
सघन, गरम पदार्थ --- बल्ब का तंतु, पिघला इस्पात, किसी तारे की सतह ---
दमकता है, और उसके प्रकाश में एक चौड़े परिसर की \emph{हर}
\omterm{def:g10:light-spectra:wavelength}{तरंगदैर्घ्य} होती है: रंगों की एक अटूट पट्टी, जिसमें कुछ भी
अनुपस्थित नहीं। ऐसा \omterm{def:g10:light-spectra:white}{वर्णक्रम} \emph{सतत वर्णक्रम}\index{सतत
वर्णक्रम} कहलाता है। उसमें पिंड की रासायनिक प्रकृति का कोई चिह्न नहीं
होता: सफ़ेद तप्त लोहे की छड़ और सफ़ेद तप्त ताँबे की छड़ एक जैसी दमकती
हैं। वह जो कूटबद्ध करता है --- और पूरी तरह करता है --- वह ताप है।
\end{definition}

\begin{definition}[परम ताप]\label{def:g10:light-spectra:kelvin}
विकिरण के नियमों के लिए \emph{परम ताप}\index{परम ताप} पानी के हिमांक से
नहीं, बल्कि संभव सबसे ठंडी अवस्था $\qty{-273}{\degreeCelsius}$ से गिना जाता है। मात्रक
\emph{केल्विन}\index{केल्विन} है (संकेत $\unit{K}$), और अंश का माप वही रहता
है:
\[
T \text{ (}\unit{K}\text{ में)} = \theta \text{ (}\unit{\degreeCelsius}\text{ में)} + 273 .
\]
इस प्रकार $\qty{20}{\degreeCelsius}$ का कमरा $\qty{293}{K}$ पर है, और सूर्य की सतह, जो $\qty{5500}{\degreeCelsius}$ के
पास है, $\qty{5800}{K}$ के पास है।
\end{definition}

\begin{proposition}[अधिक गरम यानी अधिक चमकीला और अधिक नीला (वीन का नियम)]\label{prop:g10:light-spectra:wien}
जब किसी सघन दमकते पिंड का ताप बढ़ता है, तो
\begin{enumerate}
  \item वह \emph{हर} \omterm{def:g10:light-spectra:wavelength}{तरंगदैर्घ्य} पर अधिक विकिरण करता है: पूरा
        \omterm{def:g10:light-spectra:white}{वर्णक्रम} चमक उठता है;
  \item जिस \omterm{def:g10:light-spectra:wavelength}{तरंगदैर्घ्य} $\lambda_{\max}$ पर वह सबसे प्रबल विकिरण करता है
        वह छोटी \omterm{def:g10:light-spectra:wavelength}{तरंगदैर्घ्य} (नीले) की ओर खिसकती है, इस नियम के
        अनुसार:
        \[
        \lambda_{\max} \times T = \qty{2.90e-3}{m.K},
        \]
        जहाँ $T$ केल्विनों में \omterm{def:g10:light-spectra:kelvin}{परम ताप} है।
\end{enumerate}
\end{proposition}
\admitted

\begin{remark}[यह नियम कहाँ से आता है]\label{rem:g10:light-spectra:wienorigin}
वीन का नियम उन्नीसवीं सदी के अंत में भट्ठियों और तंतुओं पर की गई सावधान
मापों का सार है। उसकी व्युत्पत्ति के लिए विकिरण का क्वांटम सिद्धांत
चाहिए --- असल में यही \omterm{def:g10:light-spectra:white}{वर्णक्रम} वह चीज़ थी जिसे चिरसम्मत भौतिकी
समझा \emph{नहीं} सकी, और यही पहेली थी जिसने प्लांक को क्वांटम गढ़ने पर
मजबूर किया। ईमानदार व्युत्पत्ति स्नातक वर्ष 3 के खंड में की गई है।
\end{remark}

\begin{omfigure}
\begin{tikzpicture}[scale=1.0, font=\small]
  \foreach \y in {0, 1.1, 2.2}
    \foreach \xa/\xb/\c in {0/0.6/violet, 0.6/1.35/blue,
        1.35/2.4/green!75!black, 2.4/2.85/yellow, 2.85/3.45/orange,
        3.45/6/red}
      \fill[\c] (\xa,\y) rectangle (\xb,\y+0.55);
  \foreach \xa/\xb/\o in {0/0.6/0.85, 0.6/1.35/0.7, 1.35/2.4/0.5,
      2.4/2.85/0.35, 2.85/3.45/0.25, 3.45/6/0.1}
    \fill[black, opacity=\o] (\xa,0) rectangle (\xb,0.55);
  \foreach \xa/\xb/\o in {0/0.6/0.5, 0.6/1.35/0.35, 1.35/2.4/0.2,
      2.4/2.85/0.12, 2.85/3.45/0.08, 3.45/6/0.08}
    \fill[black, opacity=\o] (\xa,1.1) rectangle (\xb,1.65);
  \node[left] at (-0.15, 0.275) {$\qty{3000}{K}$};
  \node[left] at (-0.15, 1.375) {$\qty{4500}{K}$};
  \node[left] at (-0.15, 2.475) {$\qty{6000}{K}$};
  \draw[-Stealth] (0,-0.5) -- (6.7,-0.5)
    node[right] {$\lambda$ (nm)};
  \foreach \x/\l in {0/400, 1.5/500, 3/600, 4.5/700, 6/800}
    \draw (\x,-0.42) -- (\x,-0.58) node[below] {\l};
\end{tikzpicture}

{\small तीन तापों पर \omterm{def:g10:light-spectra:continuous}{सतत वर्णक्रम} की दृश्य पट्टी: स्रोत जितना
ठंडा, पट्टी उतनी मंद --- और नीला सिरा सबसे पहले फीका पड़ता है।}
\end{omfigure}

\begin{example}[सूर्य का ताप लेना]\label{ex:g10:light-spectra:suntemp}
सूर्य का \omterm{def:g10:light-spectra:continuous}{सतत वर्णक्रम} $\lambda_{\max} = \qty{500}{nm}$ के पास सबसे प्रबल है। वीन का नियम
उसकी सतह का ताप देता है:
\[
T = \frac{\qty{2.90e-3}{m.K}}{\qty{5.00e-7}{m}} = \qty{5800}{K},
\]
यानी लगभग $\qty{5500}{\degreeCelsius}$ --- और यह $150$ करोड़ किलोमीटर दूर से नापा गया, एक
प्रिज़्म और एक भाग से।
\end{example}

\begin{omfigure}
\begin{tikzpicture}
\begin{axis}[omaxis, width=9.4cm, height=5.6cm,
    xlabel={$\lambda$ (nm)}, ylabel={तीव्रता (मनमाने मात्रक)},
    xmin=0, xmax=2150, ymin=0, ymax=9.5,
    xtick={400, 800, 1200, 1600, 2000}, ytick=\empty]
  \fill[gray, opacity=0.12] (axis cs:400,0) rectangle (axis cs:800,9.5);
  \node[gray, anchor=south] at (axis cs:600,8.4) {दृश्य};
  \addplot[omThm, thick, domain=180:2050, samples=160]
    {3e16 * x^(-5) / (exp(2400/x) - 1)};
  \addplot[omDef, thick, domain=180:2050, samples=160]
    {3e16 * x^(-5) / (exp(3600/x) - 1)};
  \draw[dashed, omThm] (axis cs:480,0) -- (axis cs:480,8.0);
  \draw[dashed, omDef] (axis cs:725,0) -- (axis cs:725,1.05);
  \node[omThm, anchor=south west] at (axis cs:620,6.4)
    {गरम तारा, $\qty{6000}{K}$};
  \node[omDef, anchor=south west] at (axis cs:950,1.3)
    {ठंडा तारा, $\qty{4000}{K}$};
\end{axis}
\end{tikzpicture}

{\small दो तारों के लिए \omterm{def:g10:light-spectra:wavelength}{तरंगदैर्घ्य} के सामने तीव्रता: अधिक गरम
तारा हर \omterm{def:g10:light-spectra:wavelength}{तरंगदैर्घ्य} पर अधिक विकिरण करता है, और उसका शिखर $\lambda_{\max}$
(बिंदुदार) छोटी \omterm{def:g10:light-spectra:wavelength}{तरंगदैर्घ्य} पर बैठता है।}
\end{omfigure}

\section{उत्तेजित गैसों के रेखा वर्णक्रम}

\begin{definition}[उत्सर्जन रेखा वर्णक्रम]\label{def:g10:light-spectra:emission}
कम दाब पर रखी गैस जब विद्युत विसर्जन से उत्तेजित होती है (जैसे नियॉन
नलिका या सोडियम वाली सड़क-बत्ती में), तो वह भी दमकती है --- पर
\omterm{def:g10:light-spectra:white}{वर्णक्रमदर्शी} से देखने पर उसका प्रकाश इंद्रधनुष जैसा बिलकुल नहीं
होता। केवल कुछ अलग-थलग \omterm{def:g10:light-spectra:wavelength}{तरंगदैर्घ्य} ही मौजूद होती हैं: काली
पृष्ठभूमि पर पतली चमकीली रेखाएँ, जिन्हें \emph{वर्णक्रम
रेखाएँ}\index{वर्णक्रम रेखा} कहते हैं। ऐसा \omterm{def:g10:light-spectra:white}{वर्णक्रम}
\emph{उत्सर्जन रेखा वर्णक्रम}\index{उत्सर्जन वर्णक्रम}\index{रेखा
वर्णक्रम} कहलाता है।
\end{definition}

\begin{example}[हाइड्रोजन और सोडियम]\label{ex:g10:light-spectra:hydrogen}
उत्तेजित हाइड्रोजन दृश्य परिसर में ठीक चार रेखाएँ देता है: $\qty{656}{nm}$ पर एक
लाल रेखा, $\qty{486}{nm}$ पर एक नीली-हरी रेखा, और $\qty{434}{nm}$ तथा $\qty{410}{nm}$ पर दो बैंगनी
रेखाएँ। उत्तेजित सोडियम मुख्यतः $\qty{589}{nm}$ पर एक ही प्रबल पीली-नारंगी रेखा
देता है --- इसीलिए पुरानी सड़क-बत्तियाँ पूरे मोहल्ले को उसी एक रंग में
नहला देती थीं: उनके प्रकाश में और लगभग कुछ होता ही नहीं।
\end{example}

\begin{proposition}[तत्वों की अँगुली-छाप]\label{prop:g10:light-spectra:fingerprint}
हर रासायनिक तत्व अपनी \omterm{def:g10:light-spectra:white}{वर्णक्रम} रेखाओं की एक निश्चित सूची
उत्सर्जित करता है, जो हर प्रयोगशाला में और हर युग में एक जैसी रहती है,
और किन्हीं दो तत्वों की सूची एक जैसी नहीं होती। इसलिए \omterm{def:g10:light-spectra:emission}{रेखा
वर्णक्रम} उत्सर्जक तत्व को उसी तरह पहचान देता है जैसे अँगुली-छाप किसी
व्यक्ति को।
\end{proposition}
\admitted

\begin{remark}[रेखाएँ ही क्यों?]\label{rem:g10:light-spectra:whylines}
कोई परमाणु केवल कुछ ही \omterm{def:g10:light-spectra:wavelength}{तरंगदैर्घ्य} क्यों उत्सर्जित करता है ---
और वे \omterm{def:g10:light-spectra:wavelength}{तरंगदैर्घ्य} तत्व का भेद क्यों खोल देती हैं --- यह परमाणु
के क्वांटम सिद्धांत से समझाया जाता है, इसी खंड के अंतिम वर्ष में: हर
तत्व के पास ऊर्जा स्तरों की अपनी एक सीढ़ी होती है, और हर रेखा दो डंडों
के बीच की एक छलाँग दर्ज करती है। फ़िलहाल हम अँगुली-छाप का उपयोग जासूसों
की तरह करेंगे: अँगुलियों से यह पूछे बिना कि वे बनीं कैसे।
\end{remark}

\section{अवशोषण वर्णक्रम}

\begin{definition}[अवशोषण वर्णक्रम]\label{def:g10:light-spectra:absorption}
\omterm{def:g10:light-spectra:white}{श्वेत प्रकाश} को किसी ठंडी, कम दाब वाली गैस \emph{से होकर}
भेजिए और जो निकले उसे विक्षेपित कीजिए: सतत इंद्रधनुष फिर दिखता है, पर
उस पर पतली काली रेखाएँ खिंची होती हैं। गैस ने गुज़रते प्रकाश में से कुछ
\omterm{def:g10:light-spectra:wavelength}{तरंगदैर्घ्य} हटा दी हैं। ऐसा \omterm{def:g10:light-spectra:white}{वर्णक्रम} --- सतत पृष्ठभूमि
में से कुछ रेखाएँ घटाकर --- \emph{अवशोषण वर्णक्रम}\index{अवशोषण
वर्णक्रम} कहलाता है।
\end{definition}

\begin{proposition}[उत्सर्जन और अवशोषण रेखाएँ एक ही जगह पड़ती हैं]\label{prop:g10:light-spectra:kirchhoff}
कोई गैस ठीक वही \omterm{def:g10:light-spectra:wavelength}{तरंगदैर्घ्य} अवशोषित करती है जो वह उत्सर्जित कर
सकती है: किसी तत्व के \omterm{def:g10:light-spectra:absorption}{अवशोषण वर्णक्रम} की काली रेखाएँ ठीक उन्हीं
तरंगदैर्घ्यों पर बैठती हैं जिन पर उसके \omterm{def:g10:light-spectra:emission}{उत्सर्जन वर्णक्रम}
की चमकीली रेखाएँ। इसलिए \cref{prop:g10:light-spectra:fingerprint} की अँगुली-छाप अँधेरे में भी उतनी ही अच्छी
तरह पढ़ी जा सकती है जितनी उजाले में।
\end{proposition}
\admitted

\begin{remark}[एक ही क्रियाविधि, दो चिह्न]\label{rem:g10:light-spectra:mechanism}
यह संयोग नहीं है: $\lambda$ पर अवशोषण वही ऊर्जा-डंडा चढ़ता है जिसे $\lambda$
पर उत्सर्जन उतरता है --- इस खंड के अंत के क्वांटम अध्याय इसे ठीक-ठीक
बताते हैं। 1859 में किरखोफ़ और बुन्सन द्वारा स्थापित यही बात वह कुंजी
है जो पूरा अगला अनुभाग खोल देती है।
\end{remark}

\begin{omfigure}
\begin{tikzpicture}[scale=1.0, font=\small]
  \fill[black] (0,1.5) rectangle (6,2.1);
  \draw[violet, line width=1.6pt] (0.15,1.5) -- (0.15,2.1);
  \draw[blue, line width=1.6pt]   (0.51,1.5) -- (0.51,2.1);
  \draw[cyan, line width=1.6pt]   (1.29,1.5) -- (1.29,2.1);
  \draw[red, line width=1.6pt]    (3.84,1.5) -- (3.84,2.1);
  \node[left] at (-0.15,1.8) {उत्सर्जन};
  \foreach \xa/\xb/\c in {0/0.6/violet, 0.6/1.35/blue,
      1.35/2.4/green!75!black, 2.4/2.85/yellow, 2.85/3.45/orange,
      3.45/6/red}
    \fill[\c] (\xa,0) rectangle (\xb,0.6);
  \foreach \x in {0.15, 0.51, 1.29, 3.84}
    \draw[black, line width=1.6pt] (\x,0) -- (\x,0.6);
  \node[left] at (-0.15,0.3) {अवशोषण};
  \foreach \x in {0.15, 0.51, 1.29, 3.84}
    \draw[dashed, gray] (\x,0.6) -- (\x,1.5);
  \draw[-Stealth] (0,-0.5) -- (6.7,-0.5)
    node[right] {$\lambda$ (nm)};
  \foreach \x/\l in {0/400, 1.5/500, 3/600, 4.5/700, 6/800}
    \draw (\x,-0.42) -- (\x,-0.58) node[below] {\l};
\end{tikzpicture}

{\small हाइड्रोजन के दो \omterm{def:g10:light-spectra:white}{वर्णक्रम}: चमकीली उत्सर्जन रेखाएँ (ऊपर) और काली
अवशोषण रेखाएँ (नीचे), ठीक उन्हीं तरंगदैर्घ्यों पर --- $410$,
$434$, $486$ और $\qty{656}{nm}$।}
\end{omfigure}

\begin{example}[सोडियम वाष्प की एक कुप्पी]\label{ex:g10:light-spectra:sodiumflask}
\omterm{def:g10:light-spectra:white}{श्वेत प्रकाश} ठंडी सोडियम वाष्प से भरी एक कुप्पी को पार करता है।
पारगत \omterm{def:g10:light-spectra:white}{वर्णक्रम} पूरा इंद्रधनुष है, पर $\qty{589}{nm}$ पर एक काली रेखा
के साथ --- ठीक वहीं जहाँ सोडियम लैंप सबसे चमकीला होता है। अँधेरे कमरे
में कुप्पी को बगल से देखने पर चुराया हुआ प्रकाश फिर से उत्सर्जित होता
दिखता है: उसी \omterm{def:g10:light-spectra:wavelength}{तरंगदैर्घ्य} पर एक मंद पीली-नारंगी झलक।
\end{example}

\section{वर्णक्रम विश्लेषण: एक तारा पढ़ना}

\begin{definition}[वर्णक्रम विश्लेषण]\label{def:g10:light-spectra:analysis}
\emph{वर्णक्रम विश्लेषण}\index{वर्णक्रम विश्लेषण} किसी स्रोत (लौ, लैंप
या तारे) में मौजूद रासायनिक तत्वों की पहचान है, जो उसके
\omterm{def:g10:light-spectra:white}{वर्णक्रम} की रेखाओं को तत्वों की प्रयोगशाला-सूचियों से मिलाकर की
जाती है।
\end{definition}

\begin{method}[तारकीय वर्णक्रम पढ़ना]\label{met:g10:light-spectra:reading}
तारा सघन गरम पदार्थ का एक गोला (सतह) है, जिस पर एक पतला, ठंडा वायुमंडल
लिपटा है। इसलिए उसका प्रकाश \omterm{def:g10:light-spectra:continuous}{सतत वर्णक्रम} के रूप में आता है
जिस पर अवशोषण रेखाएँ खिंची होती हैं --- और हर हिस्सा अलग-अलग पढ़ा जाता
है:
\begin{enumerate}
  \item सतत पृष्ठभूमि का शिखर $\lambda_{\max}$ ढूँढ़िए; वीन का नियम (\cref{prop:g10:light-spectra:wien})
        सतह का ताप $T = \qty{2.90e-3}{m.K} / \lambda_{\max}$ दे देता है।
  \item काली रेखाओं की \omterm{def:g10:light-spectra:wavelength}{तरंगदैर्घ्य} सूचीबद्ध कीजिए।
  \item उन्हें सूचियों से मिलाइए: किसी तत्व को उपस्थित तभी घोषित कीजिए
        जब उसकी \emph{सभी} प्रबल रेखाएँ दिखें।
  \item निष्कर्ष निकालिए: अवशोषण तारे के वायुमंडल में होता है, इसलिए
        पहचाने गए तत्व वायुमंडल के तत्व हैं।
\end{enumerate}
\end{method}

\begin{example}[सूर्य का नुस्ख़ा]\label{ex:g10:light-spectra:sunrecipe}
सौर \omterm{def:g10:light-spectra:white}{वर्णक्रम} में हज़ारों काली रेखाएँ हैं, जिनका नक्शा
फ्राउनहोफ़र ने 1814 में बनाया। सबसे प्रबल में से: $410$, $434$,
$486$ और $\qty{656}{nm}$ --- पूरी हाइड्रोजन अँगुली-छाप --- और उनके साथ
$\qty{589}{nm}$ पर सोडियम की रेखा तथा $393$ और $\qty{397}{nm}$ पर कैल्शियम का एक
युग्म। सूर्य के वायुमंडल में सबसे बढ़कर हाइड्रोजन है, साथ में उन तत्वों
के अंश जो पृथ्वी की चट्टानों से जाने-पहचाने हैं। आधुनिक मापें इस शीर्षक
को और पैना कर देती हैं: सूर्य द्रव्यमान के हिसाब से मोटे तौर पर
तीन-चौथाई हाइड्रोजन है।
\end{example}

\begin{remark}[वह तत्व जो पहले आकाश में मिला]\label{rem:g10:light-spectra:helium}
1868 के सूर्यग्रहण के समय सूर्य के वायुमंडल के \omterm{def:g10:light-spectra:white}{वर्णक्रम} में
$\qty{587.6}{nm}$ पर एक चमकीली रेखा दिखी --- जो किसी ज्ञात तत्व से मेल नहीं खाती
थी। \cref{prop:g10:light-spectra:fingerprint} के अनुसार अनजानी अँगुली-छाप का अर्थ है अनजाना तत्व, सो एक
तत्व घोषित कर दिया गया और उसका नाम यूनानी सूर्य \emph{हीलिओस} पर हीलियम
रखा गया। पृथ्वी पर वह अंततः 1895 में अलग किया जा सका --- अंतरिक्ष में
खोजा गया और घर पर बाद में मिला, ऐसा इकलौता रासायनिक तत्व।
\end{remark}

\section{अभ्यास}

\begin{exercise}[$\star$]\label{exo:g10:light-spectra:1}
हर विकिरण के लिए बताइए कि वह \omterm{def:g10:light-spectra:wavelength}{पराबैंगनी} है, दृश्य है या
\omterm{def:g10:light-spectra:wavelength}{अवरक्त}, और दृश्य हो तो उसका रंग भी दीजिए: $\qty{254}{nm}$ (पारा लैंप),
$\qty{480}{nm}$, $\qty{589}{nm}$, $\qty{656}{nm}$, $\qty{1064}{nm}$ (कटाई करने वाला लेज़र)।
\end{exercise}

\begin{exercise}[$\star$]\label{exo:g10:light-spectra:2}
धूप की एक पतली किरण काँच के प्रिज़्म पर पड़ती है।
\begin{enumerate}
  \item बताइए कि प्रिज़्म के पीछे रखे परदे पर क्या दिखेगा।
  \item पट्टी का कौन-सा सिरा सबसे अधिक विचलित हुआ है?
  \item धूप के स्थान पर हरे लेज़र सूचक की किरण रख दी जाती है। अब परदे
        पर क्या दिखता है, और इससे लेज़र प्रकाश के बारे में क्या पता
        चलता है?
\end{enumerate}
\end{exercise}

\begin{exercise}[$\star$]\label{exo:g10:light-spectra:3}
हर स्रोत को उस \omterm{def:g10:light-spectra:white}{वर्णक्रम} से मिलाइए जो वह देता है (सतत, रेखा
उत्सर्जन, या अवशोषण): (क) तापदीप्त बल्ब का तंतु; (ख) विज्ञापन की नियॉन
नलिका; (ग) ढलाईघर में पिघला लोहा; (घ) सोडियम वाली सड़क-बत्ती; (ङ)
पृथ्वी पर प्राप्त धूप।
\end{exercise}

\begin{exercise}[$\star$]\label{exo:g10:light-spectra:4}
इंद्रधनुष के बारे में:
\begin{enumerate}
  \item प्रिज़्म की भूमिका कौन निभाता है?
  \item इंद्रधनुष तभी दिखता है जब सूर्य प्रेक्षक के पीछे हो और वर्षा
        सामने। हर सामग्री क्या देती है?
  \item समझाइए कि इंद्रधनुष यह क्यों सिद्ध करता है कि धूप
        \omterm{def:g10:light-spectra:white}{श्वेत प्रकाश} है।
\end{enumerate}
\end{exercise}

\begin{exercise}[$\star$]\label{exo:g10:light-spectra:5}
एक लोहार लोहे की कील गरम करता है: वह पहले धुँधली लाल दमकती है, फिर
चमकीली नारंगी, फिर लगभग सफ़ेद।
\begin{enumerate}
  \item तीनों चरणों को ताप के क्रम में रखिए।
  \item यह क्रम \cref{prop:g10:light-spectra:wien} के कौन-से दो प्रभाव दिखाता है?
  \item ``सफ़ेद तप्त, लाल तप्त से अधिक गरम होता है'' --- इस कथन को
        उचित ठहराइए।
\end{enumerate}
\end{exercise}

\begin{exercise}[$\star\star$]\label{exo:g10:light-spectra:6}
सूर्य के \omterm{def:g10:light-spectra:continuous}{सतत वर्णक्रम} का शिखर $\lambda_{\max} = \qty{500}{nm}$ पर है।
\begin{enumerate}
  \item सूर्य की सतह का ताप केल्विनों में निकालिए।
  \item उसे डिग्री सेल्सियस में बदलिए।
  \item जाँचिए कि शिखर दृश्य परिसर के भीतर पड़ता है, और उसका रंग
        बताइए।
\end{enumerate}
\end{exercise}

\begin{exercise}[$\star\star$]\label{exo:g10:light-spectra:7}
मोमबत्ती की लौ लगभग $\qty{1800}{K}$ पर दमकती है; ``गुनगुना'' पर सेट किया गया
बिजली का चूल्हा लगभग $\qty{500}{K}$ पर रहता है।
\begin{enumerate}
  \item हर स्रोत के लिए $\lambda_{\max}$ निकालिए और उसका क्षेत्र बताइए।
  \item चूल्हा साफ़ तौर पर गरम है --- उसके ऊपर रखा हाथ विकिरण महसूस
        करता है --- फिर भी वह काला दिखता है। समझाइए।
  \item मोमबत्ती का शिखर भी दृश्य परिसर के बाहर है। फिर हमें लौ दिखती
        ही क्यों है?
\end{enumerate}
\end{exercise}

\begin{exercise}[$\star\star$]\label{exo:g10:light-spectra:8}
प्रयोगशाला की सूचियाँ दृश्य परिसर में देती हैं: हाइड्रोजन $\{410, 434, 486, 656\}\,\unit{nm}$;
सोडियम $\{589\}\,\unit{nm}$; हीलियम $\{447, 502, 588, 668\}\,\unit{nm}$। दो विसर्जन लैंपों का विश्लेषण किया
जाता है। लैंप A में $434$, $486$ और $\qty{656}{nm}$ पर रेखाएँ दिखती हैं,
साथ ही दृश्यता के किनारे के पास एक मंद बैंगनी रेखा; लैंप B में $\qty{589}{nm}$
पर एक ही प्रबल रेखा दिखती है। हर लैंप की गैस पहचानिए, और सावधानी से
उचित ठहराइए कि लैंप B में हीलियम क्यों नहीं है।
\end{exercise}

\begin{exercise}[$\star\star$]\label{exo:g10:light-spectra:9}
\omterm{def:g10:light-spectra:white}{श्वेत प्रकाश} ठंडी सोडियम वाष्प की कुप्पी से भेजा जाता है।
\begin{enumerate}
  \item पारगत प्रकाश का \omterm{def:g10:light-spectra:white}{वर्णक्रम} बताइए।
  \item उसकी काली रेखा किस \omterm{def:g10:light-spectra:wavelength}{तरंगदैर्घ्य} पर बैठती है, और ठीक
        वहीं क्यों?
  \item लैंप बंद कर दिया जाता है, कमरा अँधेरा कर दिया जाता है, और अब
        वाष्प को विसर्जन से उत्तेजित किया जाता है। नया
        \omterm{def:g10:light-spectra:white}{वर्णक्रम} बताइए।
\end{enumerate}
\end{exercise}

\begin{exercise}[$\star\star$]\label{exo:g10:light-spectra:10}
सूर्य के एक उच्च-विभेदन \omterm{def:g10:light-spectra:white}{वर्णक्रम} में $410$, $434$, $486$,
$517$, $518$ और $\qty{656}{nm}$ पर काली रेखाएँ दिखती हैं। सूचियाँ:
हाइड्रोजन $\{410, 434, 486, 656\}\,\unit{nm}$; मैग्नीशियम $\{517, 518\}\,\unit{nm}$।
\begin{enumerate}
  \item ये रेखाएँ कौन-से तत्व उजागर करती हैं?
  \item ये रेखाएँ सूर्य के किस हिस्से में बनती हैं, और वे चमकीली होने
        के बजाय काली क्यों हैं?
\end{enumerate}
\end{exercise}

\begin{exercise}[$\star\star$]\label{exo:g10:light-spectra:11}
आपके शरीर की सतह का ताप लगभग $\qty{37}{\degreeCelsius}$ होता है।
\begin{enumerate}
  \item उसे केल्विनों में बदलिए और अपने ही तापीय विकिरण का
        $\lambda_{\max}$ निकालिए।
  \item वह किस क्षेत्र में पड़ता है?
  \item इससे निकालिए कि अँधेरे कमरे में कोई दिखने भर को क्यों नहीं
        दमकता, और फिर भी तापीय कैमरा लोगों को कैसे ढूँढ़ लेता है।
\end{enumerate}
\end{exercise}

\begin{exercise}[$\star\star\star$]\label{exo:g10:light-spectra:12}
तापदीप्त बल्ब का टंगस्टन तंतु $\qty{2700}{K}$ पर चलता है।
\begin{enumerate}
  \item $\lambda_{\max}$ निकालिए और उसका क्षेत्र बताइए।
  \item समझाइए कि ऐसे बल्ब घटिया लैंप पर बढ़िया छोटे हीटर क्यों होते
        हैं।
  \item किस तंतु-ताप पर $\lambda_{\max}$ दृश्य परिसर के बीचोंबीच $\qty{550}{nm}$ पर आ
        जाएगा? टंगस्टन $\qty{3700}{K}$ पर पिघलता है; इससे निष्कर्ष निकालिए कि
        कोई भी तंतु-बल्ब दिन के उजाले की नकल क्यों नहीं कर सकता, और
        प्रकाश-व्यवस्था दूसरी तकनीकों की ओर क्यों चली गई।
\end{enumerate}
\end{exercise}

\begin{exercise}[$\star\star\star$]\label{exo:g10:light-spectra:13}
किसी तारे के \omterm{def:g10:light-spectra:white}{वर्णक्रम} में $\qty{420}{nm}$ पर शिखर वाली सतत पृष्ठभूमि
है, जिस पर $393$, $397$, $410$, $434$, $486$ और $\qty{656}{nm}$ पर
काली रेखाएँ खिंची हैं। सूचियाँ: हाइड्रोजन $\{410, 434, 486, 656\}\,\unit{nm}$; कैल्शियम $\{393, 397\}\,\unit{nm}$;
सोडियम $\{589\}\,\unit{nm}$।
\begin{enumerate}
  \item तारे की सतह का ताप निकालिए। वह सूर्य से अधिक गरम है या अधिक
        ठंडा?
  \item उसके वायुमंडल में निश्चित रूप से कौन-से तत्व हैं?
  \item क्या इस वायुमंडल में सोडियम प्रचुर है? उचित ठहराइए।
\end{enumerate}
\end{exercise}

\begin{exercise}[$\star\star\star$]\label{exo:g10:light-spectra:14}
एक विसर्जन लैंप में दो गैसों का मिश्रण है। उसके \omterm{def:g10:light-spectra:white}{वर्णक्रम} में
$410$, $434$, $447$, $486$, $502$, $588$, $656$ और $\qty{668}{nm}$
पर रेखाएँ दिखती हैं। \cref{exo:g10:light-spectra:8} की सूचियों का उपयोग करते हुए:
\begin{enumerate}
  \item दोनों गैसें पहचानिए।
  \item एक सहपाठी का दावा है कि $\qty{588}{nm}$ वाली रेखा सिद्ध करती है कि
        लैंप में सोडियम भी है (रेखा $\qty{589}{nm}$ पर), क्योंकि सस्ता
        \omterm{def:g10:light-spectra:white}{वर्णक्रमदर्शी} $\qty{1}{nm}$ की दूरी पर पड़ी
        तरंगदैर्घ्यों को अलग नहीं कर सकता। बेहतर उपकरण के
        बिना, हीलियम इस रेखा की पर्याप्त व्याख्या क्यों है --- और
        सोडियम का सवाल हमेशा के लिए कौन-सी माप तय कर देगी?
\end{enumerate}
\end{exercise}

\begin{exercise}[$\star\star\star$]\label{exo:g10:light-spectra:15}
ओरायन तारामंडल में बीटलजूस साफ़ लाल दमकता है और राइजेल नीला-सफ़ेद।
उनके सतत वर्णक्रमों के शिखर क्रमशः लगभग $\qty{850}{nm}$ और $\qty{240}{nm}$ पर हैं।
\begin{enumerate}
  \item दोनों सतहों के ताप निकालिए।
  \item किसी भी ताप को दोबारा निकाले बिना, दोनों तापों का अनुपात सीधे
        दोनों तरंगदैर्घ्यों से निकालिए।
  \item दोनों रंग समझाइए, यह देखते हुए कि \emph{किसी भी} तारे का शिखर
        दृश्य परिसर में नहीं है।
  \item दोनों तारे नंगी आँख से चमकीले दिखते हैं। दृश्य परिसर के बाहर
        शिखर होने से तारा अदृश्य क्यों नहीं हो जाता?
\end{enumerate}
\end{exercise}

\section{समस्या: तारों की रोशनी पढ़ना}

\begin{problem}[{सप्ताहांत समस्या --- तारों की रोशनी पढ़ना: वह दार्शनिक
जिसने ``कभी नहीं'' कहा था, और एक प्रिज़्म कैसे किसी तारे का ताप लेकर
उसका रासायनिक नुस्ख़ा पढ़ लेता है}]
\label{pb:g10:light-spectra:1}
1835 में दार्शनिक ऑगस्त कोंत को ऐसे ज्ञान का उदाहरण चाहिए था जो मनुष्य
की पहुँच से सदा बाहर रहे, और उन्होंने बड़ी सावधानी से चुना: तारों की
रासायनिक बनावट। कोई तारा कभी बोतल में नहीं भरेगा। 1859 में किरखोफ़ और
बुन्सन ने सूर्य की काली रेखाओं को प्रयोगशाला की लौ की चमकीली रेखाओं से
मिला दिया, और वह अगम्य ज्ञान प्रकाश के साथ ही लौट आया। यह समस्या पूरा
रास्ता दोहराती है: प्रयोगशाला की अँगुली-छापें, दमकते पिंड, सूर्य की
काली रेखाएँ --- और अंत में एक तारा, जिसका ताप और बनावट आप अपनी मेज़ से
ही तय कर देंगे।

\medskip
\noindent\textbf{भाग I --- फ़ाइल में दर्ज अँगुली-छापें।}
\begin{enumerate}
  \item दृश्य तरंगदैर्घ्यों का परिसर नैनोमीटर में दीजिए और
        दोनों सिरों का रंग बताइए। $\qty{550}{nm}$ कहाँ बैठती है?
  \item हाइड्रोजन विसर्जन लैंप में $410$, $434$, $486$ और $\qty{656}{nm}$
        पर रेखाएँ दिखती हैं। हर रेखा का रंग बताइए।
  \item सोडियम लैंप मुख्यतः $\qty{589}{nm}$ पर एक ही रेखा देता है। आँख को लैंप
        कैसा दिखता है, और ऐसी सड़क-बत्ती के नीचे खड़ी कार का रंग
        आँकना बेकार क्यों है?
  \item समझाइए कि \omterm{def:g10:light-spectra:white}{वर्णक्रम} रेखाओं का एक समूह किसी तत्व की
        पहचान क्यों करा देता है --- और किसी तत्व को उपस्थित घोषित करने
        से पहले उसकी \emph{सभी} प्रबल रेखाएँ क्यों मिलनी चाहिए।
  \item बताइए कि कौन-सा \omterm{def:g10:light-spectra:white}{वर्णक्रम} दिखेगा (क) जब
        \omterm{def:g10:light-spectra:white}{श्वेत प्रकाश} ठंडी सोडियम वाष्प की कुप्पी से गुज़रे,
        और (ख) जब वही वाष्प, अँधेरे कमरे में अकेली, विसर्जन से
        उत्तेजित की जाए।
\end{enumerate}

\medskip
\noindent\textbf{भाग II --- गरम पिंड और उनके रंग।}
\begin{enumerate}[resume]
  \item बताइए कि किसी सघन दमकते पिंड का ताप बढ़ाने से उसके
        \omterm{def:g10:light-spectra:continuous}{सतत वर्णक्रम} पर कौन-से दो प्रभाव पड़ते हैं।
  \item हैलोजन तंतु $\qty{3400}{K}$ पर चलता है। $\lambda_{\max}$ निकालिए और उसका क्षेत्र
        बताइए।
  \item सूर्य के \omterm{def:g10:light-spectra:continuous}{सतत वर्णक्रम} का शिखर $\qty{500}{nm}$ पर है। उसकी सतह
        का ताप निकालिए।
  \item वह शिखर हरे में बैठता है --- फिर भी सूर्य हरा नहीं दिखता।
        समझाइए।
  \item भट्ठी में छोड़ी गई सलाख पहले दूर से ही महसूस होती है पर दमकती
        नहीं दिखती, फिर धुँधली लाल दमकती है, फिर नारंगी-सफ़ेद। प्रश्न 6
        की मदद से पूरा क्रम समझाइए।
\end{enumerate}

\medskip
\noindent\textbf{भाग III --- सूर्य की काली रेखाएँ।}
\begin{enumerate}[resume]
  \item सौर \omterm{def:g10:light-spectra:white}{वर्णक्रम} एक सतत पट्टी है जिसे हज़ारों महीन काली
        रेखाएँ काटती हैं। सूर्य का कौन-सा हिस्सा सतत पृष्ठभूमि बनाता
        है, और कौन-सा हिस्सा रेखाएँ?
  \item सूर्य की प्रबल रेखाएँ $393$, $397$, $410$, $434$,
        $486$, $589$ और $\qty{656}{nm}$ पर हैं। सूचियाँ: हाइड्रोजन
        $\{410, 434, 486, 656\}\,\unit{nm}$; सोडियम $\{589\}\,\unit{nm}$; कैल्शियम $\{393, 397\}\,\unit{nm}$। सूर्य के वायुमंडल
        में कौन-से तत्व मौजूद हैं?
  \item हाइड्रोजन की रेखाएँ कहीं अधिक प्रबल हैं। इससे सूर्य की बनावट
        के बारे में क्या संकेत मिलता है --- और आधुनिक मापें क्या कहती
        हैं?
  \item 1868 में सूर्य के वायुमंडल में देखी गई $\qty{587.6}{nm}$ वाली रेखा किसी
        सूची से मेल नहीं खाती थी। वह तर्क फिर से गढ़िए जिसने एक नए तत्व
        की घोषणा को उचित ठहराया, और बताइए कि कहानी का अंत कैसे हुआ।
  \item सूर्य की रेखाएँ काली क्यों हैं, जबकि वही परमाणु प्रयोगशाला के
        विसर्जन में चमकीली रेखाएँ देते हैं? पूर्ण सूर्यग्रहण में, कुछ
        क्षणों के लिए, अँधेरे आकाश के सामने सूर्य का केवल पतला
        वायुमंडल दिखता रहता है --- और वही रेखाएँ \emph{चमकीली} कौंध
        जाती हैं। समझाइए।
\end{enumerate}

\medskip
\noindent\textbf{भाग IV --- मेज़ पर रखा एक तारा।}
आपके सामने चमकीले तारे वेगा का \omterm{def:g10:light-spectra:white}{वर्णक्रम} रखा है: $\qty{290}{nm}$ पर
शिखर वाली सतत पृष्ठभूमि, जिस पर $410$, $434$, $486$ और $\qty{656}{nm}$ पर
प्रबल काली रेखाएँ खिंची हैं।
\begin{enumerate}[resume]
  \item वेगा की सतह का ताप निकालिए।
  \item शिखर \omterm{def:g10:light-spectra:wavelength}{पराबैंगनी} में पड़ता है। तारे का आभासी रंग बताइए,
        और ``शिखर दृश्य के बाहर'' तथा ``चमकीला तारा'' में मेल बिठाइए।
  \item वेगा के वायुमंडल पर छाए हुए तत्व को पहचानिए।
  \item वेगा की सूर्य से तुलना कीजिए: ताप का अनुपात, रंग, और दमकती सतह
        के प्रति मात्रक चमक।
  \item समापन --- वह पहचान-पत्र भरिए जिसे कोंत ने असंभव बताया था:
        वेगा का ताप, प्रमुख तत्व और आभासी रंग, सब प्रकाश की एक ही किरण
        से पढ़े हुए। ``कभी नहीं'' कितने वर्ष टिका?
\end{enumerate}
\end{problem}
